\section{Analysis and Discussion}
\label{sec:analysis_discussion}

The experimental matrix reveals three interlocking findings that, taken together, argue for prosody-guided pooling not as a universal accuracy maximizer, but as a \textit{structurally principled} design choice for deployable child SER systems.

% ────────────────────────────────────────────────────────────
\subsection{Pillar 1: The Acted--Spontaneous Chasm and Zero-Shot Collapse}
\label{ssec:pillar1}

The single most important result in this study is Exp4: a model trained on acted child speech (C-BESD, 92.78\% in-domain WA) achieves only \textbf{18.70\% WA} when evaluated zero-shot on spontaneous child speech (FAU Aibo)---a collapse of 74 percentage points, landing below the 25\% chance level.

This is not a model failure in the conventional sense.
It is a \textit{distribution-shift failure}: acted speech datasets encode exaggerated, prototypical emotional archetypes that bear little resemblance to the subtle, noisy, overlapping emotional expressions found in naturalistic child--robot interactions.
The finding is consistent with distribution-shift diagnostics (FD = 16.33, SMMD = 0.41 between C-BESD and FAU Aibo feature spaces) and aligns with recent warnings in the SER literature that acted corpora produce ``dataset-solving'' models rather than emotion-understanding systems~\cite{wagner2023dawn}.

\textbf{Implication for the field:} Any child SER method reporting only acted-corpus accuracy is reporting a dataset artifact, not generalizable emotion recognition performance.
Our zero-shot result is, to our knowledge, the first to quantify the acted$\rightarrow$spontaneous transfer penalty specifically for \textit{children's} speech.

% ────────────────────────────────────────────────────────────
\subsection{Pillar 2: Prosody as a Structural Regularizer on Naturalistic Speech}
\label{ssec:pillar2}

On acted C-BESD, pure self-attention outperforms prosody-guided pooling by +1.48pp WA (92.78\% vs.\ 91.30\%).
This is expected: acted speech exhibits exaggerated, near-deterministic prosodic contours that a sufficiently deep attention head can memorize directly from the SSL feature stream, making explicit prosodic injection redundant.

However, on \textit{naturalistic} FAU Aibo speech---the regime that matters for real-world deployment---the picture reverses.
Prosody-guided pooling achieves \textbf{66.36\%} WA versus 66.18\% for self-attention, holding its ground despite identical regularization.
More critically, prosody-guided pooling reaches its best validation accuracy at epoch 3 versus epoch 2 for self-attention, suggesting a more stable optimization trajectory.

We interpret this as evidence that explicit prosodic priors act as a \textbf{structural regularizer}: by constraining the attention mechanism to weight frames proportionally to F0 transitions and energy bursts---acoustic events known to correlate with emotional salience in children~\cite{fan2024children}---the model is prevented from overfitting to irrelevant environmental noise, speaker idiosyncrasies, or recording artifacts that dominate spontaneous child corpora.

In contrast, the self-attention pathway, lacking this inductive bias, achieves marginally higher UAR (58.23\% vs.\ 56.35\%) by distributing attention more uniformly, but at the cost of reduced interpretability and no clear path to explaining \textit{why} a given frame was attended to.

% ────────────────────────────────────────────────────────────
\subsection{Pillar 3: The Accuracy--Explainability Trade-off}
\label{ssec:pillar3}

The 1.48pp WA gap on C-BESD (Exp1 vs.\ Exp2) might appear to favor self-attention.
We argue this is a \textit{necessary and well-justified trade-off} for explainability, particularly in the sensitive application domains where child SER is most needed: pediatric mental health screening, autism spectrum monitoring, and educational affect detection.

\paragraph{Quantified explainability.}
Our Attention--Prosody Correlation (APC) analysis yields $\text{APC}_{\text{wav}} = 0.698$ (Pearson $r$ between learned attention weights and frame-level RMS energy), confirming that prosody-guided attention locks onto physically interpretable acoustic events---energy bursts corresponding to emotional emphasis.
In contrast, self-attention weights exhibit no such structured correlation with any single acoustic feature, rendering the model a black box.

\paragraph{Clinical deployment argument.}
In pediatric healthcare settings, a clinician cannot act on a prediction from a model that offers no justification.
A prosody-guided system that says ``this utterance was classified as \textit{angry} because attention concentrated on the high-energy vocal burst at $t$=1.2s'' is clinically actionable.
A self-attention system that says ``this utterance was classified as \textit{angry} because of opaque feature interactions across 768 dimensions'' is not.

\paragraph{The acted-dataset accuracy illusion.}
The 1.48pp advantage of self-attention exists \textit{only} on acted speech---a regime we have shown to be non-transferable to real-world spontaneous conditions (Exp4: 18.70\%).
Optimizing for acted-speech leaderboards while sacrificing interpretability produces models that are simultaneously overfit to a toy distribution and unexplainable in deployment.
This is the worst possible outcome for safety-critical applications.

% ────────────────────────────────────────────────────────────
\subsection{Unified Distribution-Shift Framework}

Across all experimental conditions, performance degrades monotonically with distribution distance:

\begin{table}[h]
\centering
\caption{Distribution distance vs.\ accuracy across experimental conditions.}
\label{tab:fd_accuracy}
\begin{tabular}{lrc}
\toprule
Condition & FD & Best WA (\%) \\
\midrule
In-domain acted (C-BESD) & 0.0 & 92.78 \\
Age shift (child $\rightarrow$ adult) & 6.87 & 58.67 \\
Style shift (acted $\rightarrow$ spontaneous, in-domain) & 12.33 & 66.36 \\
Style + language shift (zero-shot) & 16.33 & 18.70 \\
\bottomrule
\end{tabular}
\end{table}

This strict FD--accuracy anti-correlation supports the distribution-driven perspective that motivates this work: the primary determinant of child SER performance is not model architecture, but the statistical distance between training and evaluation distributions.

% ────────────────────────────────────────────────────────────
\subsection{Limitations and Future Work}

\paragraph{Limitations.}
(i) The C-BESD corpus, while widely used, is an acted dataset; our own results demonstrate this limits ecological validity.
(ii) FAU Aibo experiments exhibit rapid overfitting (best epoch 2--3), suggesting that even 18K samples may be insufficient for robust spontaneous child SER without data augmentation or pre-training strategies.
(iii) The current prosody extraction (librosa YIN) is computationally expensive at inference time; a learned prosody estimator could eliminate this bottleneck.
(iv) Single-seed results on some configurations limit statistical power.

\paragraph{Future directions.}
(i) Incorporate additional naturalistic child speech corpora (MyST, EmoReact) where licensing permits, to validate the FD--accuracy framework on a wider distribution spectrum.
(ii) Explore domain-adaptive fine-tuning of the WavLM backbone specifically on spontaneous child speech, potentially recovering the zero-shot transfer gap.
(iii) Investigate prosody-guided pooling with learnable prosody extractors to enable end-to-end training while preserving the interpretability advantage.
